% ============================================================
%  CAPÍTULO 2: ARREGLO EXPERIMENTAL
% ============================================================
\chapter{Arreglo experimental}
\label{cap:arreglo}

En este capítulo se describe el arreglo experimental desarrollado
para la medición fotoacústica de líquidos. Se presentan la celda
fotoacústica, los equipos de medición involucrados, el sistema
de medición de energía del láser y los sistemas de automatización
implementados para el osciloscopio y el láser.

% ============================================================
\section{Descripción de la celda fotoacústica}
\label{sec:celda}
% ============================================================

% GUÍA:
% - Geometría de la celda (dimensiones, material, ventanas ópticas)
% - Posición del hidrófono / fotodiodo dentro de la celda
% - Fotografía o diagrama esquemático (figura obligatoria aquí)
% - Cómo se introduce y extrae la muestra líquida

% TODO: redactar contenido aquí

% Ejemplo de cómo insertar figura cuando la tengas:
% \begin{figure}[H]
%   \centering
%   \includegraphics[width=0.7\textwidth]{celda_fotoacustica}
%   \caption{Diagrama de la celda fotoacústica utilizada.}
%   \label{fig:celda}
% \end{figure}


% ============================================================
\section{Equipos de medición: osciloscopio, hidrófono y fotodiodo}
\label{sec:equipos}
% ============================================================

% GUÍA:
% - Osciloscopio Tektronix TDS5000B: especificaciones relevantes
%   (canales, ancho de banda, resolución vertical, interfaz VXI-11)
% - Hidrófono: modelo, sensibilidad, rango de frecuencia
% - Fotodiodo: modelo, respuesta espectral, función en el arreglo
% - Tabla comparativa o tabla de especificaciones técnicas

% TODO: redactar contenido aquí


% ============================================================
\section{Sistema de medición de energía del láser}
\label{sec:energia_laser}
% ============================================================

% GUÍA:
% - Modelo del láser (Nd:YAG, longitud de onda, energía por pulso, tasa de repetición)
% - Sensor de energía (joulmeter / power meter): modelo y rango
% - Cómo se mide la energía antes y/o durante el experimento
% - Esquema del arreglo óptico completo (figura)

% TODO: redactar contenido aquí


% ============================================================
\section{Automatización del osciloscopio}
\label{sec:auto_osci}
% ============================================================

% GUÍA:
% - Protocolo VXI-11 sobre Ethernet
% - PyVISA como capa de abstracción: instalación, ResourceManager, open_resource
% - Comandos SCPI principales utilizados (adquisición de forma de onda,
%   configuración de escala, trigger)
% - Diagrama de flujo del proceso de adquisición
% - Fragmento de código representativo (usa el entorno lstlisting)

% TODO: redactar contenido aquí

% Ejemplo bloque de código:
% \begin{lstlisting}[caption={Adquisición de forma de onda con PyVISA},
%                    label={lst:pyvisa_waveform}]
% import pyvisa
% rm = pyvisa.ResourceManager()
% osci = rm.open_resource("TCPIP::192.168.1.100::INSTR")
% osci.write("DATA:SOURCE CH1")
% raw = osci.query_binary_values("CURVE?", datatype='B')
% \end{lstlisting}


% ============================================================
\section{Automatización del láser}
\label{sec:auto_laser}
% ============================================================

% GUÍA:
% - DLL del fabricante: nombre, versión, arquitectura (32-bit)
% - Uso de ctypes para cargar la DLL desde Python
% - Funciones principales expuestas por la DLL (inicialización,
%   disparo, lectura de parámetros)
% - Consideraciones de compatibilidad 32-bit (proceso separado /
%   subproceso, bridge con multiprocessing o subprocess + pickle)
% - Fragmento de código representativo

% TODO: redactar contenido aquí
